\scanpage{29}

\inserted{Ceci rappelé,} Soit \deleted{\unsure{m}}Maintenant $M$ un $\Lambda$-module libre de rang $2$, et considérons l'ens.

\[
\operatorname{Rep}(M)'\simeq\operatorname{Isom}_\Lambda(\Lambda^2,M)/\underset{\deleted{\unsure{3}}}{\pm}1\tag*{(73)}
\]

qui est un torseur \inserted{à droite} sous $\mathrm{GL}(2,\Lambda)/{\pm1}$. Comme l'hom canonique

\[
\mathbb{Z}\unsure{\mathbin{-}|\longrightarrow}\Lambda
\]

définit des hom correspondants

\[
\left\{\begin{aligned}
\mathrm{GL}(2,\mathbb{Z})&\longrightarrow\deleted{\unsure{\mathrm S}}\mathrm{GL}(2,\Lambda)\\
\mathrm{GL}(\unsure{n|2},\mathbb{Z})'&\longrightarrow\deleted{\unsure{\mathrm S}}\mathrm{GL}(n,\Lambda)'=\mathrm{GL}(2,\Lambda)/\{\pm1\}
\end{aligned}\right.\tag*{(74)}
\]

on peut considérer $\operatorname{Rep}(M)'$ comme un ensemble sur lequel $\mathrm{GL}(2,\mathbb{Z})'$ opère à droite, donc à gauche si on y tient par la modification habituelle $g\longmapsto g^{-1}$.

Donc $\operatorname{Rep}(M)'$ définit une structure \inserted{$2$-polyèdrale} triangulée (pas néc. finie) qu'il faudrait expliciter en termes des constructions du \S\ I. On peut aussi, \deleted{\unsure{gardant|prenant}} \inserted{regarder} la restriction de l'opération précédente à $\mathrm{SL}(2,\mathbb{Z})'$, on trouve une structure $2$-polyèdrale triangulée orientée (qui, pour

% archive sheet number: 28
